% Methods Section 1: Data Acquisition and Preprocessing
% Based on notebook: Copy_of_drugbankDataset_wrangling_(1).ipynb

\documentclass[11pt,a4paper]{article}
\usepackage[utf8]{inputenc}
\usepackage{geometry}
\geometry{margin=1in}

\title{\textbf{Methods: Data Acquisition and Preprocessing}}
\author{}
\date{}

\begin{document}
\maketitle

\section*{Data Acquisition and Preprocessing}

Drug-drug interaction data was obtained from DrugBank (version 5.1.10), a comprehensive pharmaceutical knowledge base containing detailed drug information including molecular properties, pharmacological targets, and documented interactions. The complete DrugBank XML database was parsed using Python's lxml library to extract all pairwise drug interactions along with their textual descriptions. For each drug, Anatomical Therapeutic Chemical (ATC) classification codes were retrieved to enable therapeutic categorization. Drugs were classified as cardiovascular agents if their ATC codes began with ``C'' (Cardiovascular System) or as antithrombotic agents if their codes began with ``B01'' (Antithrombotic Agents). The complete dataset was then filtered to retain only interactions where at least one participating drug belonged to either the cardiovascular or antithrombotic therapeutic class. This filtering yielded a focused dataset of 759,774 drug-drug interactions involving 4,314 unique drugs, providing a clinically relevant subset for subsequent severity recalibration and knowledge graph construction while maintaining sufficient scale and diversity to support robust machine learning-based classification.

\end{document}
